% !PRINTED 314
\weberpage{314}
%<<CONT
ferner die durch (1) des vorigen Paragraphen definirte Grösse $A$; sodann aber auch die symmetrischen Functionen der sieben Grössen $\Psi^2$, und ferner, wenn wir nun der Einfachheit halber die Annahme hinzufügen, auf die wir noch zurückkommen, dass von den Grössen $\Psi$ keine verschwinde, die symmetrischen Functionen der Grössen $v$; aber nicht nur die symmetrischen Functionen der $v$, sondern alle Functionen der $v$, die die Permutationen der Gruppe $T$ gestatten.

Die sieben Grössen

\weq{2}{\begin{aligned}
v_{1,0,0} &= v_1, & v_{0,1,0} &= v_2, & v_{0,0,1} &= v_3, \\
v_{0,1,1} &= v_4, & v_{1,0,1} &= v_5, & v_{1,1,0} &= v_6, & v_{1,1,1} &= v_7
\end{aligned}}

sind alsdann die Wurzeln einer rationalen Gleichung $7^{\mathrm{ten}}$ Grades mit der Gruppe $T$, und durch Addition der Gleichung (1) und der Gleichungen [\S.\,80, (1)]:

\[
\Sigma (-1)^{\Sigma z\xi} X_{z_1, z_2, z_3} = \Psi_{\xi_1, \xi_2, \xi_3}
\]

ergiebt sich mit Rücksicht auf (9) des vorigen Paragraphen:

\weq{3}{8\,X_{0,0,0} = A\Bigl\{\begin{aligned}[t]
&\sqrt{v_1}\sqrt{v_5}\sqrt{v_6}\sqrt{v_7} + \sqrt{v_2}\sqrt{v_4}\sqrt{v_6}\sqrt{v_7} \\
&+ \sqrt{v_3}\sqrt{v_4}\sqrt{v_5}\sqrt{v_7} + \sqrt{v_2}\sqrt{v_3}\sqrt{v_5}\sqrt{v_6} \\
&+ \sqrt{v_1}\sqrt{v_3}\sqrt{v_4}\sqrt{v_6} + \sqrt{v_1}\sqrt{v_2}\sqrt{v_4}\sqrt{v_5} \\
&+ \sqrt{v_1}\sqrt{v_2}\sqrt{v_3}\sqrt{v_7}
\end{aligned}\Bigr\} + B.}

Es giebt 15 Vorzeichenänderungen der sieben Quadratwurzeln $\sqrt{v}$, bei denen dieser Ausdruck ungeändert bleibt, nämlich, wenn alle sieben Vorzeichen gleichzeitig geändert werden oder wenn die Vorzeichen eines Tripels oder eines Quadrupels geändert werden. \gesperrt{Folglich erhält der Ausdruck (3) nur acht verschiedene Werthe, wie man auch die sieben darin vorkommenden Quadratwurzeln bestimmen mag.}

\begin{center}
\S.\,83.

\gesperrt{Metacyklische Gleichungen $8^{\mathrm{ten}}$ Grades.}
\end{center}

Um nun nach diesen Vorbereitungen die primitiven metacyklischen Gleichungen $8^{\mathrm{ten}}$ Grades zu finden, haben wir zunächst die metacyklischen Theiler der Gruppe $S$ zu ermitteln. Die Gruppe $S$ selbst ist, wie schon erwähnt, nicht metacyklisch. Ihre
%CONT>>
